\section{Evaluación de EvoDiff: calidad, diversidad y validación funcional}

\subsection{Metodología de evaluación}

Evaluar un modelo de generación de proteínas es mucho más complejo que evaluar un modelo de generación de imágenes. Amini enfatiza que es necesario diseñar la estrategia de evaluación con cuidado desde \textbf{múltiples niveles}:

\begin{table}[H]
\centering
\caption{Marco de evaluación multinivel para modelos de generación de proteínas}
\begin{tabular}{@{}lll@{}}
\toprule
\textbf{Dimensión de evaluación} & \textbf{Pregunta central} & \textbf{Método de evaluación} \\
\midrule
Calidad individual & ¿Es la proteína individual biológicamente viable? & Predicción de estructura + self-consistency \\
Cobertura de la distribución & ¿Cubren las muestras generadas la distribución de datos? & Visualización y comparación de la distribución en el espacio de características \\
Realización funcional & ¿Puede la proteína generada cumplir la función buscada? & Validación experimental en laboratorio húmedo \\
\bottomrule
\end{tabular}
\end{table}

\begin{warningbox}{No basta con ver la "exactitud"}
La evaluación de un modelo de generación de proteínas \textbf{no puede reducirse a un solo indicador de exactitud}. Un modelo puede generar proteínas individuales de alta calidad, pero si todas las proteínas generadas se concentran en una región estrecha del espacio funcional, su valor como herramienta de diseño se ve muy disminuido. Por eso, la \textbf{diversidad a nivel de distribución} y la \textbf{capacidad de validación a nivel funcional} son tan importantes como la calidad individual.
\end{warningbox}

\subsection{Evaluación de la calidad individual: Self-Consistency}

EvoDiff adopta un método ingenioso de evaluación por self-consistency:

\begin{enumerate}
    \item Se genera con EvoDiff una nueva secuencia de proteína $\mathbf{s}_{\text{gen}}$
    \item Se predice con una herramienta de predicción de estructura (como AlphaFold) la estructura tridimensional $\mathcal{S}$ de $\mathbf{s}_{\text{gen}}$
    \item A partir de la estructura $\mathcal{S}$, se infiere una secuencia $\mathbf{s}_{\text{inv}}$ capaz de plegarse en esa estructura
    \item Se calcula la consistencia entre $\mathbf{s}_{\text{gen}}$ y $\mathbf{s}_{\text{inv}}$
\end{enumerate}

Si la consistencia es alta, indica que la secuencia generada puede plegarse de manera estable en una estructura tridimensional determinada, es decir, es \textbf{estructuralmente viable} (structurally realistic and consistent). Los resultados muestran que las proteínas generadas por EvoDiff tienen un buen desempeño en consistencia estructural.

\subsection{Evaluación de la cobertura de la distribución: comparación con proteínas naturales}

Para evaluar la diversidad de las muestras generadas por EvoDiff, el equipo realizó un análisis a nivel de distribución:

\begin{enumerate}
    \item Se toma el conjunto de prueba de secuencias de proteínas naturales y se extraen vectores de características con un modelo preentrenado
    \item Los vectores de características se proyectan a un espacio 2D mediante un método de reducción de dimensionalidad (como t-SNE o UMAP), obteniendo la distribución de las proteínas naturales
    \item Se realiza la misma operación con una gran cantidad de proteínas generadas por EvoDiff (miles de secuencias)
    \item Se superponen ambas distribuciones para evaluar el grado de cobertura
\end{enumerate}

\begin{importantbox}{Ventaja de EvoDiff en la cobertura de la distribución}
La comparación con varios métodos de referencia muestra una ventaja notable de EvoDiff:
\begin{itemize}
    \item \textbf{frente a modelos de Next Token Prediction}: ambos alcanzan una cobertura aproximadamente equivalente, pero el modelo de next token es ligeramente superior
    \item \textbf{frente a modelos de generación en un solo paso con Masked}: la cobertura de EvoDiff es claramente más amplia, lo que indica que el denoising iterativo es superior a la generación en un solo paso
    \item \textbf{frente a métodos puramente estructurales} (como RFdiffusion, entre otros): los métodos estructurales muestran un \textbf{sesgo grave}, con una gran cantidad de muestras concentradas en una región estrecha del espacio funcional
\end{itemize}
\end{importantbox}

El sesgo grave de los métodos puramente estructurales tiene su origen en la distribución sesgada de los datos estructurales: entre las aproximadamente 300,000 estructuras de proteínas resueltas experimentalmente, las proteínas globulares y las ricas en $\alpha$-helix están sobrerrepresentadas. Los modelos entrenados con estos datos naturalmente sobremuestrean (oversample) este tipo de proteínas. EvoDiff, en cambio, se entrena con datos de secuencia, cuyo volumen es mayor y cuya cobertura es más amplia, por lo que captura mejor la diversidad del espacio funcional de las proteínas naturales.

\subsection{Validación experimental: de lo computacional al laboratorio}

La evaluación de referencia definitiva es la \textbf{validación experimental en laboratorio húmedo}. El equipo tomó proteínas candidatas de entre los resultados generados por EvoDiff y realizó los siguientes experimentos:

\textbf{Validación de la estabilidad estructural}:
\begin{itemize}
    \item Se sometieron secuencias de proteínas diseñadas por completo de manera artificial a \textbf{expresión biológica} (síntesis de la proteína dentro de la célula)
    \item Se midió la estabilidad estructural de la proteína
    \item Resultado: se expresaron y validaron con éxito 4 diseños de proteínas completamente nuevas y estructuralmente estables
\end{itemize}

\begin{knowledgebox}{El ciclo completo, del diseño computacional a la validación experimental}
El valor final del diseño de proteínas con IA debe demostrarse mediante validación experimental. Esto implica un proceso complejo: (1) el modelo computacional genera secuencias candidatas; (2) mediante síntesis génica se obtiene el ADN correspondiente; (3) la proteína se expresa en un sistema celular; (4) se purifica la proteína; (5) se valida su estructura y función con métodos fisicoquímicos (como dicroísmo circular, cristalografía de rayos X, etc.). El equipo de EvoDiff completó todo este ciclo, demostrando la utilidad práctica del modelo.
\end{knowledgebox}

\subsection{Resumen de la sección}

La evaluación de EvoDiff abarcó tres niveles: calidad individual, diversidad de la distribución y validación experimental. En cuanto a la cobertura de la distribución, EvoDiff muestra una ventaja clara frente a los métodos puramente estructurales, lo cual se debe a la ventaja intrínseca de los datos de secuencia en escala y diversidad. Lo más importante es que las proteínas generadas por EvoDiff se expresaron y se validaron estructuralmente con éxito en el laboratorio, lo que demuestra la utilidad práctica del modelo.

